%-----------------------------------------------------------------------------%
\chapter{\babDua}
\label{bab:2}
%-----------------------------------------------------------------------------%
Bab ini akan menjelaskan isi dari pelaksanaan kerja praktik, mulai dari latar
belakang pekerjaan, deskripsi pekerjaan, tinjauan pustaka, metodologi,
teknologi yang digunakan, hasil pekerjaan, aspek non teknis, hingga analisis
pelaksanaan kerja praktik.



%-----------------------------------------------------------------------------%
\section{Latar Belakang Pekerjaan}
\label{sec:latarBelakangPekerjaan}
%-----------------------------------------------------------------------------%
PT TASPEN (Persero) memiliki aplikasi internal bernama Easy Mobile yang
berfungsi untuk menunjang aktivitas karyawan sehari-hari. Aplikasi ini
memungkinkan karyawan untuk melakukan berbagai pengajuan secara digital,
seperti pengajuan cuti dan Surat Perintah Tugas (SPT). Dengan adanya aplikasi
ini, proses administratif menjadi lebih efisien karena tidak lagi memerlukan
dokumen fisik dan dapat dilakukan kapan saja melalui perangkat mobile.

Sebagai perusahaan yang terus berkembang, TASPEN membutuhkan pengembangan
fitur-fitur baru pada aplikasi Easy Mobile untuk meningkatkan produktivitas
karyawan. Oleh karena itu, pelaksana diberikan tugas untuk mengembangkan
fitur Request SPT dan Request Cuti pada aplikasi Easy Mobile.



%-----------------------------------------------------------------------------%
\section{Deskripsi Pekerjaan}
\label{sec:deskripsiPekerjaan}
%-----------------------------------------------------------------------------%
Selama pelaksanaan kerja praktik, pelaksana berperan sebagai Mobile Developer
Intern dengan tanggung jawab utama mengembangkan fitur-fitur pada aplikasi
Easy Mobile. Secara spesifik, pekerjaan yang dilakukan meliputi:

\begin{enumerate}
    \item Mengembangkan fitur Request SPT (Surat Perintah Tugas) yang
    memungkinkan karyawan untuk mengajukan surat perintah tugas secara digital.
    \item Mengembangkan fitur Request Cuti yang memungkinkan karyawan untuk
    mengajukan cuti melalui aplikasi mobile.
    \item Mengintegrasikan fitur-fitur tersebut dengan API \f{backend} yang
    sudah ada.
    \item Memastikan tampilan dan fungsionalitas fitur sesuai dengan desain
    yang telah ditentukan di Figma.
\end{enumerate}



%-----------------------------------------------------------------------------%
\section{Tinjauan Pustaka}
\label{sec:tinjauanPustaka}
%-----------------------------------------------------------------------------%
Subbab ini akan menjelaskan teori-teori dan teknologi yang digunakan dalam
pelaksanaan kerja praktik.

%-----------------------------------------------------------------------------%
\subsection{\f{Mobile Application Development}}
\label{sec:mobileAppDev}
%-----------------------------------------------------------------------------%
\f{Mobile application development} adalah proses pembuatan aplikasi perangkat
lunak yang berjalan pada perangkat mobile \cite{azure:mobile}. Proses ini
melibatkan penulisan kode program, desain antarmuka pengguna, pengujian, dan
deployment aplikasi ke toko aplikasi seperti Google Play Store atau Apple App
Store. Pengembangan aplikasi mobile dapat dilakukan dengan berbagai pendekatan,
seperti pengembangan native, hybrid, atau cross-platform.

%-----------------------------------------------------------------------------%
\subsection{Flutter}
\label{sec:flutter}
%-----------------------------------------------------------------------------%
Flutter adalah \f{framework} \f{open-source} yang dikembangkan oleh Google untuk
membangun aplikasi \f{cross-platform} dari satu \f{codebase} \cite{flutter:docs}.
Dengan Flutter, pengembang dapat membuat aplikasi untuk Android, iOS, web, dan
desktop dengan menggunakan bahasa pemrograman Dart. Flutter menggunakan konsep
widget untuk membangun antarmuka pengguna, di mana setiap elemen visual adalah
widget yang dapat dikustomisasi.

%-----------------------------------------------------------------------------%
\subsection{Figma}
\label{sec:figma}
%-----------------------------------------------------------------------------%
Figma adalah alat desain berbasis web yang memungkinkan kolaborasi secara
\f{real-time} antar desainer \cite{figma:wired}. Figma digunakan untuk membuat
desain antarmuka pengguna (\f{User Interface}) dan prototipe interaktif. Dalam
konteks pengembangan aplikasi, Figma menjadi referensi utama bagi pengembang
untuk mengimplementasikan desain yang telah dibuat oleh tim desainer.

%-----------------------------------------------------------------------------%
\subsection{Git}
\label{sec:git}
%-----------------------------------------------------------------------------%
Git adalah sistem kontrol versi terdistribusi yang digunakan untuk melacak
perubahan pada kode sumber selama pengembangan perangkat lunak \cite{git:progit}.
Dengan Git, pengembang dapat bekerja secara kolaboratif, membuat cabang
(\f{branch}) untuk fitur baru, dan menggabungkan (\f{merge}) perubahan dari
berbagai kontributor. GitLab digunakan sebagai platform untuk menyimpan
repositori Git dalam pelaksanaan kerja praktik ini.

%-----------------------------------------------------------------------------%
\subsection{Postman}
\label{sec:postman}
%-----------------------------------------------------------------------------%
Postman adalah platform untuk membangun dan menggunakan API (\f{Application
Programming Interface}) \cite{postman:blog}. Postman memungkinkan pengembang
untuk menguji endpoint API, melihat respons dari server, dan mendokumentasikan
API. Dalam pelaksanaan kerja praktik ini, Postman digunakan untuk memahami
struktur API yang disediakan oleh tim \f{backend} dan menguji integrasi
antara aplikasi mobile dengan API.

%-----------------------------------------------------------------------------%
\subsection{Android Studio}
\label{sec:androidStudio}
%-----------------------------------------------------------------------------%
Android Studio adalah \f{Integrated Development Environment} (IDE) resmi untuk
pengembangan aplikasi Android \cite{android:studio}. Meskipun Flutter memiliki
dukungan untuk berbagai IDE, Android Studio menyediakan fitur-fitur lengkap
seperti emulator Android, \f{debugging tools}, dan integrasi dengan Flutter
\f{plugin}. Android Studio digunakan sebagai IDE utama dalam pengembangan
aplikasi Easy Mobile.



%-----------------------------------------------------------------------------%
\section{Metodologi}
\label{sec:metodologi}
%-----------------------------------------------------------------------------%
Pengembangan fitur pada aplikasi Easy Mobile mengikuti metodologi Agile dengan
pendekatan Scrum. Tim melakukan \f{sprint} dengan durasi dua minggu, di mana
setiap \f{sprint} dimulai dengan \f{sprint planning} dan diakhiri dengan
\f{sprint review} dan \f{retrospective}.

Proses pengembangan fitur dimulai dengan analisis kebutuhan berdasarkan desain
di Figma dan dokumentasi API. Setelah memahami kebutuhan, pelaksana melakukan
implementasi kode dengan menerapkan \f{clean architecture} untuk memastikan
kode yang terstruktur dan mudah di-\f{maintain}. Setelah implementasi selesai,
dilakukan pengujian untuk memastikan fitur berjalan sesuai dengan yang
diharapkan.



%-----------------------------------------------------------------------------%
\section{Teknologi}
\label{sec:teknologi}
%-----------------------------------------------------------------------------%
Berikut adalah teknologi yang digunakan dalam pelaksanaan kerja praktik:

\begin{enumerate}
    \item \bo{Flutter} - \f{Framework} untuk pengembangan aplikasi mobile
    \f{cross-platform}.
    \item \bo{Dart} - Bahasa pemrograman yang digunakan oleh Flutter.
    \item \bo{Figma} - Alat desain untuk referensi tampilan aplikasi.
    \item \bo{GitLab} - Platform untuk menyimpan dan mengelola kode sumber.
    \item \bo{Postman} - Alat untuk menguji dan memahami API.
    \item \bo{Android Studio} - IDE untuk pengembangan aplikasi.
\end{enumerate}



%-----------------------------------------------------------------------------%
\section{Hasil Pekerjaan}
\label{sec:hasilPekerjaan}
%-----------------------------------------------------------------------------%
Selama pelaksanaan kerja praktik, pelaksana berhasil mengembangkan dua fitur
utama pada aplikasi Easy Mobile:

\subsection{Fitur Request SPT}
Fitur Request SPT memungkinkan karyawan untuk mengajukan Surat Perintah Tugas
secara digital melalui aplikasi mobile. Fitur ini mencakup:
\begin{itemize}
    \item Form pengisian data SPT seperti tujuan, tanggal, dan keperluan tugas.
    \item Validasi input untuk memastikan data yang diisi lengkap dan benar.
    \item Integrasi dengan API untuk mengirim data ke server.
    \item Tampilan status pengajuan SPT.
\end{itemize}

% TODO: Tambahkan gambar screenshot fitur Request SPT
% \begin{figure}[htbp]
% 	\centering
%     \includegraphics[width=0.5\textwidth]{assets/pics/fitur_spt.png}
% 	\caption{Tampilan Fitur Request SPT}
% 	\label{fig:fitur_spt}
% \end{figure}

\subsection{Fitur Request Cuti}
Fitur Request Cuti memungkinkan karyawan untuk mengajukan cuti melalui
aplikasi mobile. Fitur ini mencakup:
\begin{itemize}
    \item Pemilihan jenis cuti (cuti tahunan, cuti sakit, dll.).
    \item Pemilihan tanggal mulai dan selesai cuti.
    \item Form pengisian alasan cuti.
    \item Tampilan sisa jatah cuti karyawan.
    \item Integrasi dengan API untuk mengirim pengajuan ke server.
\end{itemize}

% TODO: Tambahkan gambar screenshot fitur Request Cuti
% \begin{figure}[htbp]
% 	\centering
%     \includegraphics[width=0.5\textwidth]{assets/pics/fitur_cuti.png}
% 	\caption{Tampilan Fitur Request Cuti}
% 	\label{fig:fitur_cuti}
% \end{figure}



%-----------------------------------------------------------------------------%
\section{Aspek Non Teknis}
\label{sec:aspekNonTeknis}
%-----------------------------------------------------------------------------%
Selain keterampilan teknis, pelaksanaan kerja praktik ini juga melatih
berbagai keterampilan non teknis yang penting dalam dunia kerja profesional:

\begin{enumerate}
    \item \bo{Komunikasi} - Pelaksana belajar berkomunikasi secara efektif
    dengan mentor dan anggota tim lainnya untuk memahami kebutuhan dan
    menyampaikan progress pekerjaan.
    
    \item \bo{Kerja Sama Tim} - Pelaksana bekerja dalam tim pengembang dan
    belajar berkolaborasi dengan berbagai pihak seperti desainer dan tim
    \f{backend}.
    
    \item \bo{Adaptasi} - Pelaksana belajar beradaptasi dengan lingkungan
    kerja profesional dan budaya perusahaan.
    
    \item \bo{Pengelolaan Waktu} - Pelaksana belajar mengatur waktu untuk
    menyelesaikan tugas sesuai dengan \f{deadline} yang ditentukan.
\end{enumerate}



%-----------------------------------------------------------------------------%
\section{Analisis Pelaksanaan Kerja Praktik}
\label{sec:analisisPelaksanaan}
%-----------------------------------------------------------------------------%
Subbab ini akan menjelaskan analisis terhadap pelaksanaan kerja praktik,
termasuk kendala yang dihadapi dan cara menanganinya, penilaian terhadap
tempat kerja praktik, serta relevansi dengan perkuliahan di Fasilkom UI.

%-----------------------------------------------------------------------------%
\subsection{Ulasan Tentang Kendala dan Cara Menanganinya}
\label{sec:kendala}
%-----------------------------------------------------------------------------%
Selama pelaksanaan kerja praktik, pelaksana menghadapi beberapa kendala:

\begin{enumerate}
    \item \bo{Keterbatasan Dokumentasi API} \\
    Dokumentasi API yang tersedia tidak lengkap sehingga pelaksana kesulitan
    memahami struktur dan parameter yang diperlukan. Kendala ini diatasi
    dengan melakukan komunikasi intensif dengan mentor dan tim \f{backend},
    serta memanfaatkan Postman collection yang sudah ada untuk memahami
    struktur API.
    
    \item \bo{Kerumitan User Flow} \\
    Beberapa fitur memiliki alur pengguna (\f{user flow}) yang kompleks
    sehingga memerlukan pemahaman mendalam terhadap proses bisnis. Kendala
    ini diatasi dengan melakukan analisis mandiri terhadap desain di Figma
    dan berdiskusi dengan mentor untuk mendapatkan klarifikasi.
    
    \item \bo{Adopsi Clean Architecture} \\
    Proyek menggunakan \f{clean architecture} yang memerlukan pemahaman
    struktur kode yang berbeda dari yang biasa dipelajari di perkuliahan.
    Kendala ini diatasi dengan mempelajari kode yang sudah ada dan bertanya
    kepada mentor mengenai praktik terbaik yang digunakan.
\end{enumerate}

%-----------------------------------------------------------------------------%
\subsection{Penilaian Individu Terhadap Tempat KP}
\label{sec:penilaianTempatKp}
%-----------------------------------------------------------------------------%
PT TASPEN (Persero) merupakan tempat kerja praktik yang baik untuk
mengembangkan keterampilan di bidang \f{mobile development}. Beberapa hal
positif yang dapat diambil dari tempat kerja praktik ini:

\begin{enumerate}
    \item Mentor yang responsif dan bersedia membantu ketika pelaksana
    menghadapi kesulitan.
    \item Tim yang kolaboratif dan mendukung proses pembelajaran pelaksana.
    \item Penggunaan teknologi modern seperti Flutter dan \f{clean architecture}.
    \item Lingkungan kerja yang kondusif untuk belajar dan berkembang.
\end{enumerate}

%-----------------------------------------------------------------------------%
\subsection{Relevansi dengan Perkuliahan di Fasilkom UI}
\label{sec:relevansiPerkuliahan}
%-----------------------------------------------------------------------------%
Pelaksanaan kerja praktik ini memiliki relevansi dengan beberapa mata kuliah
yang telah dipelajari di Fasilkom UI:

\begin{enumerate}
    \item \bo{Dasar-Dasar Pemrograman (DDP)} - Konsep dasar pemrograman yang
    dipelajari menjadi fondasi dalam menulis kode aplikasi.
    
    \item \bo{Pemrograman Berbasis Platform (PBP)} - Mata kuliah ini
    memberikan dasar pengembangan aplikasi mobile yang langsung diterapkan
    dalam kerja praktik.
    
    \item \bo{Basis Data} - Pemahaman tentang struktur data dan query
    membantu dalam memahami struktur API dan respons dari server.
    
    \item \bo{Rekayasa Perangkat Lunak (RPL)} - Metodologi pengembangan
    perangkat lunak seperti Agile dan Scrum yang dipelajari di mata kuliah
    ini diterapkan dalam pelaksanaan kerja praktik.
\end{enumerate}

Dengan dukungan tim dan mentor yang baik, pelaksana dapat menyelesaikan
tugas-tugas yang diberikan dan mendapatkan pemahaman yang lebih dalam
mengenai proses pengembangan aplikasi. Secara keseluruhan, pelaksanaan
kerja praktik ini memberikan wawasan yang luas tentang pengembangan
perangkat lunak dan mampu mempersiapkan pelaksana untuk menghadapi
tantangan di dunia kerja yang sesungguhnya.
